\documentclass[11pt]{article}
% jugeo-common.tex — shared preamble for all Judgment Geometry papers
% Include via: % jugeo-common.tex — shared preamble for all Judgment Geometry papers
% Include via: % jugeo-common.tex — shared preamble for all Judgment Geometry papers
% Include via: \input{jugeo-common}

% ── Packages ────────────────────────────────────────────────────────────
\usepackage{amsmath,amssymb,amsthm,mathtools}
\usepackage{tikz-cd}
\usepackage{listings}
\usepackage{xcolor}
\usepackage{xspace}
\usepackage{booktabs}
\usepackage{multirow}
\usepackage{enumitem}
\usepackage[expansion=false]{microtype}
\usepackage{hyperref}
\usepackage{cleveref}
% \usepackage{thmtools}  % disabled: not available in this TeX installation
\usepackage{float}
\usepackage{caption}
\usepackage{subcaption}
\usepackage{tabularx}
\usepackage{stmaryrd}       % \llbracket, \rrbracket
\usepackage{bbm}            % \mathbbm{1}

% ── Page geometry ───────────────────────────────────────────────────────
\usepackage[margin=1in]{geometry}

% ── Theorem environments ────────────────────────────────────────────────
\theoremstyle{plain}
\newtheorem{theorem}{Theorem}[section]
\newtheorem{lemma}[theorem]{Lemma}
\newtheorem{proposition}[theorem]{Proposition}
\newtheorem{corollary}[theorem]{Corollary}

\theoremstyle{definition}
\newtheorem{definition}[theorem]{Definition}
\newtheorem{example}[theorem]{Example}
\newtheorem{construction}[theorem]{Construction}

\theoremstyle{remark}
\newtheorem{remark}[theorem]{Remark}
\newtheorem{notation}[theorem]{Notation}

% ── Python listings style ──────────────────────────────────────────────
\definecolor{jugeoBlue}{HTML}{2563EB}
\definecolor{jugeoGreen}{HTML}{059669}
\definecolor{jugeoRed}{HTML}{DC2626}
\definecolor{jugeoPurple}{HTML}{7C3AED}
\definecolor{jugeoGray}{HTML}{6B7280}
\definecolor{jugeoBg}{HTML}{F8FAFC}

\lstdefinestyle{jugeo-python}{
  language=Python,
  basicstyle=\ttfamily\small,
  keywordstyle=\color{jugeoBlue}\bfseries,
  stringstyle=\color{jugeoGreen},
  commentstyle=\color{jugeoGray}\itshape,
  numberstyle=\tiny\color{jugeoGray},
  backgroundcolor=\color{jugeoBg},
  numbers=left,
  numbersep=6pt,
  frame=single,
  framerule=0.4pt,
  rulecolor=\color{jugeoGray!40},
  breaklines=true,
  breakatwhitespace=true,
  showstringspaces=false,
  tabsize=4,
  xleftmargin=1.5em,
  framexleftmargin=1.5em,
  morekeywords={dataclass,frozen,field,Enum,IntEnum,auto,Any,
                Mapping,Sequence,Optional,match,case},
  literate={→}{{$\to$}}1 {←}{{$\leftarrow$}}1
           {≤}{{$\leq$}}1 {≥}{{$\geq$}}1 {≠}{{$\neq$}}1
           {∈}{{$\in$}}1 {∀}{{$\forall$}}1 {∃}{{$\exists$}}1
           {⊕}{{$\oplus$}}1 {⊖}{{$\ominus$}}1,
  escapeinside={(*@}{@*)},
}
\lstset{style=jugeo-python}

\lstdefinestyle{jugeo-output}{
  basicstyle=\ttfamily\small,
  backgroundcolor=\color{jugeoBg},
  frame=single,
  framerule=0.4pt,
  rulecolor=\color{jugeoGray!40},
  breaklines=true,
  showstringspaces=false,
  xleftmargin=1.5em,
  framexleftmargin=0.5em,
  numbers=none,
}

% ── JuGeo-specific macros ──────────────────────────────────────────────

% Judgment 8-tuple
\newcommand{\J}{\mathcal{J}}
\newcommand{\Jud}[8]{(#1,\,#2,\,#3,\,#4,\,#5,\,#6,\,#7,\,#8)}
\newcommand{\JudShort}[1]{\J_{#1}}

% Components
\newcommand{\coord}{\mathsf{c}}            % coordinate
\newcommand{\prop}{\varphi}                 % proposition
\newcommand{\carrier}{\mathcal{A}}          % carrier type
\newcommand{\evid}{\mathcal{E}}             % evidence bundle
\newcommand{\oblig}{\mathcal{O}}            % obligations
\newcommand{\obstr}{\mathcal{B}}            % obstructions
\newcommand{\trust}{\mathcal{T}}            % trust annotation
\newcommand{\prov}{\Pi}                     % provenance

% Trust algebra
\newcommand{\Talg}{\mathfrak{T}}            % trust ordered algebra
\newcommand{\Eadm}{\mathcal{E}_{\mathrm{adm}}}
\newcommand{\tleq}{\preceq}                 % trust ordering
\newcommand{\tjoin}{\oplus}                 % conservative join
\newcommand{\tatten}{\ominus}               % attenuation
\newcommand{\tpromote}{\uparrow_\pi}        % promotion
\newcommand{\tdemote}{\downarrow_\chi}       % demotion

% Trust tiers
\newcommand{\Tcontra}{\ensuremath{\mathsf{CONTRADICTED}}}
\newcommand{\Tunver}{\ensuremath{\mathsf{UNVERIFIED}}}
\newcommand{\Tcopilot}{\ensuremath{\mathsf{COPILOT}}}
\newcommand{\Toracle}{\ensuremath{\mathsf{ORACLE}}}
\newcommand{\Truntime}{\ensuremath{\mathsf{RUNTIME}}}
\newcommand{\Tsolver}{\ensuremath{\mathsf{SOLVER}}}
\newcommand{\Tproof}{\ensuremath{\mathsf{PROOF}}}

% Site and geometry
\newcommand{\Site}{\mathcal{S}}             % semantic site
\newcommand{\Cov}{\mathrm{Cov}}            % covering family
\newcommand{\Ob}{\mathrm{Ob}}              % objects
\newcommand{\Mor}{\mathrm{Mor}}            % morphisms
\newcommand{\res}{\mathrm{res}}            % restriction
\newcommand{\Cech}{\check{C}}              % Čech complex
\newcommand{\Hcoh}[1]{H^{#1}}             % cohomology group

% Descent
\newcommand{\descent}{\mathrm{desc}}
\newcommand{\glue}{\mathrm{glue}}
\newcommand{\overlap}[2]{#1 \cap #2}

% Semantic moves
\newcommand{\VERIFY}{\textsc{Verify}}
\newcommand{\CONSTRUCT}{\textsc{Construct}}
\newcommand{\NORMALIZE}{\textsc{Normalize}}
\newcommand{\GLUE}{\textsc{Glue}}
\newcommand{\RESTRICT}{\textsc{Restrict}}
\newcommand{\TRANSPORT}{\textsc{Transport}}
\newcommand{\REPAIR}{\textsc{Repair}}
\newcommand{\PROMOTE}{\textsc{Promote}}
\newcommand{\CHALLENGE}{\textsc{Challenge}}
\newcommand{\ESCALATE}{\textsc{Escalate}}

% Fragments
\newcommand{\QFLIA}{\texttt{QF\_LIA}}
\newcommand{\QFLRA}{\texttt{QF\_LRA}}
\newcommand{\QFBV}{\texttt{QF\_BV}}
\newcommand{\QFUF}{\texttt{QF\_UF}}

% Misc
\newcommand{\Python}{\textsc{Python}}
\newcommand{\JuGeo}{\textsc{JuGeo}}
\newcommand{\LEAN}{\textsc{Lean}\,4}
\newcommand{\Fstar}{F$^{\star}$}
\newcommand{\Coq}{\textsc{Coq}}
\newcommand{\Dafny}{\textsc{Dafny}}

% Common shortcuts
\newcommand{\ie}{i.e.\@\xspace}
\newcommand{\eg}{e.g.\@\xspace}
\newcommand{\cf}{cf.\@\xspace}
\newcommand{\wrt}{w.r.t.\@\xspace}

% Evaluation shortcuts
\newcommand{\numcases}{300}
\newcommand{\numspec}{100}
\newcommand{\numeq}{100}
\newcommand{\numbug}{100}
\newcommand{\medtime}{6.5\,ms}
\newcommand{\totaltime}{1.949\,s}


% ── Packages ────────────────────────────────────────────────────────────
\usepackage{amsmath,amssymb,amsthm,mathtools}
\usepackage{tikz-cd}
\usepackage{listings}
\usepackage{xcolor}
\usepackage{xspace}
\usepackage{booktabs}
\usepackage{multirow}
\usepackage{enumitem}
\usepackage[expansion=false]{microtype}
\usepackage{hyperref}
\usepackage{cleveref}
% \usepackage{thmtools}  % disabled: not available in this TeX installation
\usepackage{float}
\usepackage{caption}
\usepackage{subcaption}
\usepackage{tabularx}
\usepackage{stmaryrd}       % \llbracket, \rrbracket
\usepackage{bbm}            % \mathbbm{1}

% ── Page geometry ───────────────────────────────────────────────────────
\usepackage[margin=1in]{geometry}

% ── Theorem environments ────────────────────────────────────────────────
\theoremstyle{plain}
\newtheorem{theorem}{Theorem}[section]
\newtheorem{lemma}[theorem]{Lemma}
\newtheorem{proposition}[theorem]{Proposition}
\newtheorem{corollary}[theorem]{Corollary}

\theoremstyle{definition}
\newtheorem{definition}[theorem]{Definition}
\newtheorem{example}[theorem]{Example}
\newtheorem{construction}[theorem]{Construction}

\theoremstyle{remark}
\newtheorem{remark}[theorem]{Remark}
\newtheorem{notation}[theorem]{Notation}

% ── Python listings style ──────────────────────────────────────────────
\definecolor{jugeoBlue}{HTML}{2563EB}
\definecolor{jugeoGreen}{HTML}{059669}
\definecolor{jugeoRed}{HTML}{DC2626}
\definecolor{jugeoPurple}{HTML}{7C3AED}
\definecolor{jugeoGray}{HTML}{6B7280}
\definecolor{jugeoBg}{HTML}{F8FAFC}

\lstdefinestyle{jugeo-python}{
  language=Python,
  basicstyle=\ttfamily\small,
  keywordstyle=\color{jugeoBlue}\bfseries,
  stringstyle=\color{jugeoGreen},
  commentstyle=\color{jugeoGray}\itshape,
  numberstyle=\tiny\color{jugeoGray},
  backgroundcolor=\color{jugeoBg},
  numbers=left,
  numbersep=6pt,
  frame=single,
  framerule=0.4pt,
  rulecolor=\color{jugeoGray!40},
  breaklines=true,
  breakatwhitespace=true,
  showstringspaces=false,
  tabsize=4,
  xleftmargin=1.5em,
  framexleftmargin=1.5em,
  morekeywords={dataclass,frozen,field,Enum,IntEnum,auto,Any,
                Mapping,Sequence,Optional,match,case},
  literate={→}{{$\to$}}1 {←}{{$\leftarrow$}}1
           {≤}{{$\leq$}}1 {≥}{{$\geq$}}1 {≠}{{$\neq$}}1
           {∈}{{$\in$}}1 {∀}{{$\forall$}}1 {∃}{{$\exists$}}1
           {⊕}{{$\oplus$}}1 {⊖}{{$\ominus$}}1,
  escapeinside={(*@}{@*)},
}
\lstset{style=jugeo-python}

\lstdefinestyle{jugeo-output}{
  basicstyle=\ttfamily\small,
  backgroundcolor=\color{jugeoBg},
  frame=single,
  framerule=0.4pt,
  rulecolor=\color{jugeoGray!40},
  breaklines=true,
  showstringspaces=false,
  xleftmargin=1.5em,
  framexleftmargin=0.5em,
  numbers=none,
}

% ── JuGeo-specific macros ──────────────────────────────────────────────

% Judgment 8-tuple
\newcommand{\J}{\mathcal{J}}
\newcommand{\Jud}[8]{(#1,\,#2,\,#3,\,#4,\,#5,\,#6,\,#7,\,#8)}
\newcommand{\JudShort}[1]{\J_{#1}}

% Components
\newcommand{\coord}{\mathsf{c}}            % coordinate
\newcommand{\prop}{\varphi}                 % proposition
\newcommand{\carrier}{\mathcal{A}}          % carrier type
\newcommand{\evid}{\mathcal{E}}             % evidence bundle
\newcommand{\oblig}{\mathcal{O}}            % obligations
\newcommand{\obstr}{\mathcal{B}}            % obstructions
\newcommand{\trust}{\mathcal{T}}            % trust annotation
\newcommand{\prov}{\Pi}                     % provenance

% Trust algebra
\newcommand{\Talg}{\mathfrak{T}}            % trust ordered algebra
\newcommand{\Eadm}{\mathcal{E}_{\mathrm{adm}}}
\newcommand{\tleq}{\preceq}                 % trust ordering
\newcommand{\tjoin}{\oplus}                 % conservative join
\newcommand{\tatten}{\ominus}               % attenuation
\newcommand{\tpromote}{\uparrow_\pi}        % promotion
\newcommand{\tdemote}{\downarrow_\chi}       % demotion

% Trust tiers
\newcommand{\Tcontra}{\ensuremath{\mathsf{CONTRADICTED}}}
\newcommand{\Tunver}{\ensuremath{\mathsf{UNVERIFIED}}}
\newcommand{\Tcopilot}{\ensuremath{\mathsf{COPILOT}}}
\newcommand{\Toracle}{\ensuremath{\mathsf{ORACLE}}}
\newcommand{\Truntime}{\ensuremath{\mathsf{RUNTIME}}}
\newcommand{\Tsolver}{\ensuremath{\mathsf{SOLVER}}}
\newcommand{\Tproof}{\ensuremath{\mathsf{PROOF}}}

% Site and geometry
\newcommand{\Site}{\mathcal{S}}             % semantic site
\newcommand{\Cov}{\mathrm{Cov}}            % covering family
\newcommand{\Ob}{\mathrm{Ob}}              % objects
\newcommand{\Mor}{\mathrm{Mor}}            % morphisms
\newcommand{\res}{\mathrm{res}}            % restriction
\newcommand{\Cech}{\check{C}}              % Čech complex
\newcommand{\Hcoh}[1]{H^{#1}}             % cohomology group

% Descent
\newcommand{\descent}{\mathrm{desc}}
\newcommand{\glue}{\mathrm{glue}}
\newcommand{\overlap}[2]{#1 \cap #2}

% Semantic moves
\newcommand{\VERIFY}{\textsc{Verify}}
\newcommand{\CONSTRUCT}{\textsc{Construct}}
\newcommand{\NORMALIZE}{\textsc{Normalize}}
\newcommand{\GLUE}{\textsc{Glue}}
\newcommand{\RESTRICT}{\textsc{Restrict}}
\newcommand{\TRANSPORT}{\textsc{Transport}}
\newcommand{\REPAIR}{\textsc{Repair}}
\newcommand{\PROMOTE}{\textsc{Promote}}
\newcommand{\CHALLENGE}{\textsc{Challenge}}
\newcommand{\ESCALATE}{\textsc{Escalate}}

% Fragments
\newcommand{\QFLIA}{\texttt{QF\_LIA}}
\newcommand{\QFLRA}{\texttt{QF\_LRA}}
\newcommand{\QFBV}{\texttt{QF\_BV}}
\newcommand{\QFUF}{\texttt{QF\_UF}}

% Misc
\newcommand{\Python}{\textsc{Python}}
\newcommand{\JuGeo}{\textsc{JuGeo}}
\newcommand{\LEAN}{\textsc{Lean}\,4}
\newcommand{\Fstar}{F$^{\star}$}
\newcommand{\Coq}{\textsc{Coq}}
\newcommand{\Dafny}{\textsc{Dafny}}

% Common shortcuts
\newcommand{\ie}{i.e.\@\xspace}
\newcommand{\eg}{e.g.\@\xspace}
\newcommand{\cf}{cf.\@\xspace}
\newcommand{\wrt}{w.r.t.\@\xspace}

% Evaluation shortcuts
\newcommand{\numcases}{300}
\newcommand{\numspec}{100}
\newcommand{\numeq}{100}
\newcommand{\numbug}{100}
\newcommand{\medtime}{6.5\,ms}
\newcommand{\totaltime}{1.949\,s}


% ── Packages ────────────────────────────────────────────────────────────
\usepackage{amsmath,amssymb,amsthm,mathtools}
\usepackage{tikz-cd}
\usepackage{listings}
\usepackage{xcolor}
\usepackage{xspace}
\usepackage{booktabs}
\usepackage{multirow}
\usepackage{enumitem}
\usepackage[expansion=false]{microtype}
\usepackage{hyperref}
\usepackage{cleveref}
% \usepackage{thmtools}  % disabled: not available in this TeX installation
\usepackage{float}
\usepackage{caption}
\usepackage{subcaption}
\usepackage{tabularx}
\usepackage{stmaryrd}       % \llbracket, \rrbracket
\usepackage{bbm}            % \mathbbm{1}

% ── Page geometry ───────────────────────────────────────────────────────
\usepackage[margin=1in]{geometry}

% ── Theorem environments ────────────────────────────────────────────────
\theoremstyle{plain}
\newtheorem{theorem}{Theorem}[section]
\newtheorem{lemma}[theorem]{Lemma}
\newtheorem{proposition}[theorem]{Proposition}
\newtheorem{corollary}[theorem]{Corollary}

\theoremstyle{definition}
\newtheorem{definition}[theorem]{Definition}
\newtheorem{example}[theorem]{Example}
\newtheorem{construction}[theorem]{Construction}

\theoremstyle{remark}
\newtheorem{remark}[theorem]{Remark}
\newtheorem{notation}[theorem]{Notation}

% ── Python listings style ──────────────────────────────────────────────
\definecolor{jugeoBlue}{HTML}{2563EB}
\definecolor{jugeoGreen}{HTML}{059669}
\definecolor{jugeoRed}{HTML}{DC2626}
\definecolor{jugeoPurple}{HTML}{7C3AED}
\definecolor{jugeoGray}{HTML}{6B7280}
\definecolor{jugeoBg}{HTML}{F8FAFC}

\lstdefinestyle{jugeo-python}{
  language=Python,
  basicstyle=\ttfamily\small,
  keywordstyle=\color{jugeoBlue}\bfseries,
  stringstyle=\color{jugeoGreen},
  commentstyle=\color{jugeoGray}\itshape,
  numberstyle=\tiny\color{jugeoGray},
  backgroundcolor=\color{jugeoBg},
  numbers=left,
  numbersep=6pt,
  frame=single,
  framerule=0.4pt,
  rulecolor=\color{jugeoGray!40},
  breaklines=true,
  breakatwhitespace=true,
  showstringspaces=false,
  tabsize=4,
  xleftmargin=1.5em,
  framexleftmargin=1.5em,
  morekeywords={dataclass,frozen,field,Enum,IntEnum,auto,Any,
                Mapping,Sequence,Optional,match,case},
  literate={→}{{$\to$}}1 {←}{{$\leftarrow$}}1
           {≤}{{$\leq$}}1 {≥}{{$\geq$}}1 {≠}{{$\neq$}}1
           {∈}{{$\in$}}1 {∀}{{$\forall$}}1 {∃}{{$\exists$}}1
           {⊕}{{$\oplus$}}1 {⊖}{{$\ominus$}}1,
  escapeinside={(*@}{@*)},
}
\lstset{style=jugeo-python}

\lstdefinestyle{jugeo-output}{
  basicstyle=\ttfamily\small,
  backgroundcolor=\color{jugeoBg},
  frame=single,
  framerule=0.4pt,
  rulecolor=\color{jugeoGray!40},
  breaklines=true,
  showstringspaces=false,
  xleftmargin=1.5em,
  framexleftmargin=0.5em,
  numbers=none,
}

% ── JuGeo-specific macros ──────────────────────────────────────────────

% Judgment 8-tuple
\newcommand{\J}{\mathcal{J}}
\newcommand{\Jud}[8]{(#1,\,#2,\,#3,\,#4,\,#5,\,#6,\,#7,\,#8)}
\newcommand{\JudShort}[1]{\J_{#1}}

% Components
\newcommand{\coord}{\mathsf{c}}            % coordinate
\newcommand{\prop}{\varphi}                 % proposition
\newcommand{\carrier}{\mathcal{A}}          % carrier type
\newcommand{\evid}{\mathcal{E}}             % evidence bundle
\newcommand{\oblig}{\mathcal{O}}            % obligations
\newcommand{\obstr}{\mathcal{B}}            % obstructions
\newcommand{\trust}{\mathcal{T}}            % trust annotation
\newcommand{\prov}{\Pi}                     % provenance

% Trust algebra
\newcommand{\Talg}{\mathfrak{T}}            % trust ordered algebra
\newcommand{\Eadm}{\mathcal{E}_{\mathrm{adm}}}
\newcommand{\tleq}{\preceq}                 % trust ordering
\newcommand{\tjoin}{\oplus}                 % conservative join
\newcommand{\tatten}{\ominus}               % attenuation
\newcommand{\tpromote}{\uparrow_\pi}        % promotion
\newcommand{\tdemote}{\downarrow_\chi}       % demotion

% Trust tiers
\newcommand{\Tcontra}{\ensuremath{\mathsf{CONTRADICTED}}}
\newcommand{\Tunver}{\ensuremath{\mathsf{UNVERIFIED}}}
\newcommand{\Tcopilot}{\ensuremath{\mathsf{COPILOT}}}
\newcommand{\Toracle}{\ensuremath{\mathsf{ORACLE}}}
\newcommand{\Truntime}{\ensuremath{\mathsf{RUNTIME}}}
\newcommand{\Tsolver}{\ensuremath{\mathsf{SOLVER}}}
\newcommand{\Tproof}{\ensuremath{\mathsf{PROOF}}}

% Site and geometry
\newcommand{\Site}{\mathcal{S}}             % semantic site
\newcommand{\Cov}{\mathrm{Cov}}            % covering family
\newcommand{\Ob}{\mathrm{Ob}}              % objects
\newcommand{\Mor}{\mathrm{Mor}}            % morphisms
\newcommand{\res}{\mathrm{res}}            % restriction
\newcommand{\Cech}{\check{C}}              % Čech complex
\newcommand{\Hcoh}[1]{H^{#1}}             % cohomology group

% Descent
\newcommand{\descent}{\mathrm{desc}}
\newcommand{\glue}{\mathrm{glue}}
\newcommand{\overlap}[2]{#1 \cap #2}

% Semantic moves
\newcommand{\VERIFY}{\textsc{Verify}}
\newcommand{\CONSTRUCT}{\textsc{Construct}}
\newcommand{\NORMALIZE}{\textsc{Normalize}}
\newcommand{\GLUE}{\textsc{Glue}}
\newcommand{\RESTRICT}{\textsc{Restrict}}
\newcommand{\TRANSPORT}{\textsc{Transport}}
\newcommand{\REPAIR}{\textsc{Repair}}
\newcommand{\PROMOTE}{\textsc{Promote}}
\newcommand{\CHALLENGE}{\textsc{Challenge}}
\newcommand{\ESCALATE}{\textsc{Escalate}}

% Fragments
\newcommand{\QFLIA}{\texttt{QF\_LIA}}
\newcommand{\QFLRA}{\texttt{QF\_LRA}}
\newcommand{\QFBV}{\texttt{QF\_BV}}
\newcommand{\QFUF}{\texttt{QF\_UF}}

% Misc
\newcommand{\Python}{\textsc{Python}}
\newcommand{\JuGeo}{\textsc{JuGeo}}
\newcommand{\LEAN}{\textsc{Lean}\,4}
\newcommand{\Fstar}{F$^{\star}$}
\newcommand{\Coq}{\textsc{Coq}}
\newcommand{\Dafny}{\textsc{Dafny}}

% Common shortcuts
\newcommand{\ie}{i.e.\@\xspace}
\newcommand{\eg}{e.g.\@\xspace}
\newcommand{\cf}{cf.\@\xspace}
\newcommand{\wrt}{w.r.t.\@\xspace}

% Evaluation shortcuts
\newcommand{\numcases}{300}
\newcommand{\numspec}{100}
\newcommand{\numeq}{100}
\newcommand{\numbug}{100}
\newcommand{\medtime}{6.5\,ms}
\newcommand{\totaltime}{1.949\,s}

% experiment-data.tex — AUTO-GENERATED from real jugeo CLI runs
% DO NOT EDIT — regenerate with: python3 experiments/run_universal_metrics.py
% Generated from 30 programs across 6 domains

\newcommand{\expTotalPrograms}{30}
\newcommand{\expVerified}{30}
\newcommand{\expAccuracy}{100.0\%}
\newcommand{\expCoordsMin}{2}
\newcommand{\expCoordsMax}{10}
\newcommand{\expCoordsMean}{5.2}
\newcommand{\expCoordsMedian}{5.5}
\newcommand{\expPropsMin}{4}
\newcommand{\expPropsMax}{20}
\newcommand{\expPropsMean}{10.4}
\newcommand{\expPropsSum}{311}
\newcommand{\expPropsOkSum}{310}
\newcommand{\expTimeMin}{3.506\,s}
\newcommand{\expTimeMax}{6.714\,s}
\newcommand{\expTimeMean}{4.64\,s}
\newcommand{\expTimeMedian}{4.508\,s}
\newcommand{\expTimeTotal}{139.204\,s}
\newcommand{\expObstructionTotal}{0}
\newcommand{\expHOneZero}{30}
\newcommand{\expDomainCount}{6}
\newcommand{\expTrustCOPILOTSUGGESTED}{30}
\newcommand{\expDomainDsCount}{5}
\newcommand{\expDomainDsVerified}{5}
\newcommand{\expDomainDsAvgCoords}{7.6}
\newcommand{\expDomainDsAvgProps}{15.2}
\newcommand{\expDomainMathCount}{5}
\newcommand{\expDomainMathVerified}{5}
\newcommand{\expDomainMathAvgCoords}{4.8}
\newcommand{\expDomainMathAvgProps}{9.4}
\newcommand{\expDomainSmCount}{5}
\newcommand{\expDomainSmVerified}{5}
\newcommand{\expDomainSmAvgCoords}{7.6}
\newcommand{\expDomainSmAvgProps}{15.2}
\newcommand{\expDomainSortCount}{5}
\newcommand{\expDomainSortVerified}{5}
\newcommand{\expDomainSortAvgCoords}{2.4}
\newcommand{\expDomainSortAvgProps}{4.8}
\newcommand{\expDomainStrCount}{5}
\newcommand{\expDomainStrVerified}{5}
\newcommand{\expDomainStrAvgCoords}{3.2}
\newcommand{\expDomainStrAvgProps}{6.4}
\newcommand{\expDomainWebCount}{5}
\newcommand{\expDomainWebVerified}{5}
\newcommand{\expDomainWebAvgCoords}{5.6}
\newcommand{\expDomainWebAvgProps}{11.2}

% subsystem-data.tex — AUTO-GENERATED from real jugeo subsystem experiments
% DO NOT EDIT — regenerate with: python3 experiments/run_subsystem_experiments.py

\newcommand{\subCliTotal}{10}
\newcommand{\subCliVerified}{9}
\newcommand{\subCliAccuracy}{90.0\%}
\newcommand{\subCliMeanTime}{4.132\,s}
\newcommand{\subCliMinTime}{3.472\,s}
\newcommand{\subCliMaxTime}{5.372\,s}
\newcommand{\subCliTotalCoords}{54}
\newcommand{\subCliTotalProps}{107}
\newcommand{\subCliTotalPropsOk}{106}
\newcommand{\subSiteCalls}{90}
\newcommand{\subSiteSuccessful}{69}
\newcommand{\subSiteSuccessRate}{76.7\%}
\newcommand{\subSiteMeanTime}{0.0001\,s}
\newcommand{\subSiteTotalTime}{0.005\,s}
\newcommand{\subSpecificationsatisfactionSmallTime}{0.0003\,s}
\newcommand{\subSpecificationsatisfactionMediumTime}{0.0001\,s}
\newcommand{\subSpecificationsatisfactionLargeTime}{0.0001\,s}
\newcommand{\subInhabitantfleetSmallTime}{0.0\,s}
\newcommand{\subInhabitantfleetMediumTime}{0.0\,s}
\newcommand{\subInhabitantfleetLargeTime}{0.0\,s}
\newcommand{\subTheoremecologySmallTime}{0.0\,s}
\newcommand{\subTheoremecologyMediumTime}{0.0\,s}
\newcommand{\subTheoremecologyLargeTime}{0.0\,s}
\newcommand{\subStatespaceexplorationSmallTime}{0.0\,s}
\newcommand{\subStatespaceexplorationMediumTime}{0.0\,s}
\newcommand{\subStatespaceexplorationLargeTime}{0.0\,s}
\newcommand{\subRepairsemanticsSmallTime}{0.0\,s}
\newcommand{\subRepairsemanticsMediumTime}{0.0\,s}
\newcommand{\subRepairsemanticsLargeTime}{0.0\,s}
\newcommand{\subReplaygluingSmallTime}{0.0\,s}
\newcommand{\subReplaygluingMediumTime}{0.0\,s}
\newcommand{\subReplaygluingLargeTime}{0.0\,s}
\newcommand{\subSemanticfuturesSmallTime}{0.0\,s}
\newcommand{\subSemanticfuturesMediumTime}{0.0\,s}
\newcommand{\subSemanticfuturesLargeTime}{0.0\,s}
\newcommand{\subHypercovertreatySmallTime}{0.0\,s}
\newcommand{\subHypercovertreatyMediumTime}{0.0\,s}
\newcommand{\subHypercovertreatyLargeTime}{0.0\,s}
\newcommand{\subEncodeforsolverSmallTime}{0.0026\,s}
\newcommand{\subEncodeforsolverMediumTime}{0.0002\,s}
\newcommand{\subEncodeforsolverLargeTime}{0.0001\,s}
\newcommand{\subInterfaceroutingSmallTime}{0.0\,s}
\newcommand{\subInterfaceroutingMediumTime}{0.0\,s}
\newcommand{\subInterfaceroutingLargeTime}{0.0\,s}
\newcommand{\subOrchestrateverificationSmallTime}{0.0002\,s}
\newcommand{\subOrchestrateverificationMediumTime}{0.0\,s}
\newcommand{\subOrchestrateverificationLargeTime}{0.0\,s}
\newcommand{\subRunfulldescentSmallTime}{0.0\,s}
\newcommand{\subRunfulldescentMediumTime}{0.0\,s}
\newcommand{\subRunfulldescentLargeTime}{0.0\,s}
\newcommand{\subKernellifecycleSmallTime}{0.0\,s}
\newcommand{\subKernellifecycleMediumTime}{0.0\,s}
\newcommand{\subKernellifecycleLargeTime}{0.0\,s}
\newcommand{\subDiscoverypipelineSmallTime}{0.0\,s}
\newcommand{\subDiscoverypipelineMediumTime}{0.0\,s}
\newcommand{\subDiscoverypipelineLargeTime}{0.0\,s}
\newcommand{\subAnalogytransportSmallTime}{0.0\,s}
\newcommand{\subAnalogytransportMediumTime}{0.0\,s}
\newcommand{\subAnalogytransportLargeTime}{0.0\,s}
\newcommand{\subFormalcoresiteSmallTime}{0.0001\,s}
\newcommand{\subFormalcoresiteMediumTime}{0.0\,s}
\newcommand{\subFormalcoresiteLargeTime}{0.0\,s}
\newcommand{\subProblematlasSmallTime}{0.0\,s}
\newcommand{\subProblematlasMediumTime}{0.0\,s}
\newcommand{\subProblematlasLargeTime}{0.0\,s}
\newcommand{\subPublicalignmentSmallTime}{0.0\,s}
\newcommand{\subPublicalignmentMediumTime}{0.0\,s}
\newcommand{\subPublicalignmentLargeTime}{0.0\,s}
\newcommand{\subRegimebootstrappingSmallTime}{0.0\,s}
\newcommand{\subRegimebootstrappingMediumTime}{0.0\,s}
\newcommand{\subRegimebootstrappingLargeTime}{0.0\,s}
\newcommand{\subRelationalrefinementSmallTime}{0.0\,s}
\newcommand{\subRelationalrefinementMediumTime}{0.0\,s}
\newcommand{\subRelationalrefinementLargeTime}{0.0\,s}
\newcommand{\subSemanticclosureSmallTime}{0.0\,s}
\newcommand{\subSemanticclosureMediumTime}{0.0\,s}
\newcommand{\subSemanticclosureLargeTime}{0.0\,s}
\newcommand{\subTheoremeconomicsSmallTime}{0.0001\,s}
\newcommand{\subTheoremeconomicsMediumTime}{0.0\,s}
\newcommand{\subTheoremeconomicsLargeTime}{0.0\,s}
\newcommand{\subBenchmarksuiteSmallTime}{0.0\,s}
\newcommand{\subBenchmarksuiteMediumTime}{0.0\,s}
\newcommand{\subBenchmarksuiteLargeTime}{0.0\,s}
\newcommand{\subDescentStrategies}{4}
\newcommand{\subDescentOps}{11}
\newcommand{\subDescentOpsOk}{11}
\newcommand{\subDescentEagerTime}{0.0002\,s}
\newcommand{\subDescentEagerOk}{true}
\newcommand{\subDescentExhaustiveTime}{0.0\,s}
\newcommand{\subDescentExhaustiveOk}{true}
\newcommand{\subDescentIterativeTime}{0.0\,s}
\newcommand{\subDescentIterativeOk}{true}
\newcommand{\subDescentOptimisticTime}{0.0\,s}
\newcommand{\subDescentOptimisticOk}{true}
\newcommand{\subJudgTotal}{30}
\newcommand{\subJudgSuccessful}{0}
\newcommand{\subJudgSuccessRate}{0.0\%}
\newcommand{\subJudgMeanTimeUs}{7.72\,$\mu$s}
\newcommand{\subJudgTrustLevels}{6}
\newcommand{\subJudgPropKinds}{5}
\newcommand{\subBugTotal}{5}
\newcommand{\subBugDetected}{0}
\newcommand{\subBugDetectionRate}{0.0\%}
\newcommand{\subBugFalsePositiveRate}{0.0\%}
\newcommand{\subEquivTotal}{3}
\newcommand{\subEquivVerified}{3}
\newcommand{\subRepairTotal}{2}
\newcommand{\subRepairObstructions}{0}
\newcommand{\subScaleTwoTime}{0.0001\,s}
\newcommand{\subScaleFourTime}{0.0\,s}
\newcommand{\subScaleEightTime}{0.0\,s}
\newcommand{\subScaleTwelveTime}{0.0\,s}
\newcommand{\subScaleSixteenTime}{0.0\,s}
\newcommand{\subScaleTwentyTime}{0.0\,s}
\newcommand{\subTotalExperimentTime}{95.6\,s}

% comprehensive-data.tex — AUTO-GENERATED from real jugeo experiments
% DO NOT EDIT — regenerate with: python3 experiments/run_comprehensive_experiments.py
% Generated: 2026-03-23 09:19:06

% --- Size scaling metrics ---
\newcommand{\compTotalPrograms}{18}
\newcommand{\compTotalVerified}{18}
\newcommand{\compOverallAccuracy}{100.0\%}
\newcommand{\compTinyCount}{5}
\newcommand{\compTinyVerified}{5}
\newcommand{\compTinyMeanCoords}{3}
\newcommand{\compTinyMeanProps}{6}
\newcommand{\compTinyMeanTime}{4.95\,s}
\newcommand{\compTinyMeanLines}{5}
\newcommand{\compTinyMinTime}{4.45\,s}
\newcommand{\compTinyMaxTime}{5.51\,s}
\newcommand{\compSmallCount}{5}
\newcommand{\compSmallVerified}{5}
\newcommand{\compSmallMeanCoords}{3.6}
\newcommand{\compSmallMeanProps}{7.2}
\newcommand{\compSmallMeanTime}{4.85\,s}
\newcommand{\compSmallMeanLines}{16.0}
\newcommand{\compSmallMinTime}{4.30\,s}
\newcommand{\compSmallMaxTime}{5.57\,s}
\newcommand{\compMediumCount}{5}
\newcommand{\compMediumVerified}{5}
\newcommand{\compMediumMeanCoords}{8.6}
\newcommand{\compMediumMeanProps}{17.2}
\newcommand{\compMediumMeanTime}{4.47\,s}
\newcommand{\compMediumMeanLines}{45.0}
\newcommand{\compMediumMinTime}{4.26\,s}
\newcommand{\compMediumMaxTime}{4.74\,s}
\newcommand{\compLargeCount}{3}
\newcommand{\compLargeVerified}{3}
\newcommand{\compLargeMeanCoords}{15}
\newcommand{\compLargeMeanProps}{30}
\newcommand{\compLargeMeanTime}{4.31\,s}
\newcommand{\compLargeMeanLines}{79.0}
\newcommand{\compLargeMinTime}{4.27\,s}
\newcommand{\compLargeMaxTime}{4.38\,s}
% --- Aggregate timing ---
\newcommand{\compTimeMean}{4.68\,s}
\newcommand{\compTimeMedian}{4.50\,s}
\newcommand{\compTimeMin}{4.265\,s}
\newcommand{\compTimeMax}{5.571\,s}
\newcommand{\compTimeTotal}{84.3\,s}
\newcommand{\compTimeStdev}{0.45\,s}
% --- Aggregate structure ---
\newcommand{\compCoordsMean}{6.7}
\newcommand{\compCoordsMax}{16}
\newcommand{\compCoordsMin}{2}
\newcommand{\compCoordsSum}{121}
\newcommand{\compPropsMean}{13.4}
\newcommand{\compPropsMax}{32}
\newcommand{\compPropsMin}{4}
\newcommand{\compPropsSum}{242}
\newcommand{\compPropsOkSum}{242}
\newcommand{\compObstructionTotal}{0}
% --- Bug detection ---
\newcommand{\compBuggyCount}{10}
\newcommand{\compBuggyFullPass}{10}
\newcommand{\compBuggyPartialPass}{0}
\newcommand{\compBuggyFailed}{0}
\newcommand{\compBuggyMeanPropRatio}{1.00}
\newcommand{\compCorrectMeanPropRatio}{1.00}
\newcommand{\compBuggyMeanTime}{5.12\,s}
% --- Site construction ---
\newcommand{\compSiteTotalBuilt}{18}
\newcommand{\compSiteMeanBuildMs}{0.05}
\newcommand{\compSiteMaxBuildMs}{0.14}
\newcommand{\compSiteMeanCoords}{6.5}
\newcommand{\compSiteMeanMorphisms}{14.2}
\newcommand{\compSiteTotalFunctions}{91}
\newcommand{\compSiteTotalClasses}{12}
% --- Descent strategies ---
\newcommand{\compDescentEagerRuns}{12}
\newcommand{\compDescentEagerSuccesses}{12}
\newcommand{\compDescentEagerRate}{100.0\%}
\newcommand{\compDescentEagerMeanMs}{0.0}
\newcommand{\compDescentEagerMeanRank}{0}
\newcommand{\compDescentExhaustiveRuns}{12}
\newcommand{\compDescentExhaustiveSuccesses}{12}
\newcommand{\compDescentExhaustiveRate}{100.0\%}
\newcommand{\compDescentExhaustiveMeanMs}{0.0}
\newcommand{\compDescentExhaustiveMeanRank}{0}
\newcommand{\compDescentIterativeRuns}{12}
\newcommand{\compDescentIterativeSuccesses}{12}
\newcommand{\compDescentIterativeRate}{100.0\%}
\newcommand{\compDescentIterativeMeanMs}{0.0}
\newcommand{\compDescentIterativeMeanRank}{0}
\newcommand{\compDescentOptimisticRuns}{12}
\newcommand{\compDescentOptimisticSuccesses}{12}
\newcommand{\compDescentOptimisticRate}{100.0\%}
\newcommand{\compDescentOptimisticMeanMs}{0.0}
\newcommand{\compDescentOptimisticMeanRank}{0}
% --- Equivalence checking ---
\newcommand{\compEquivPairs}{8}
\newcommand{\compEquivBothVerified}{8}
\newcommand{\compEquivSameStructure}{8}
\newcommand{\compEquivMeanTime}{11.06\,s}
% --- Trust distribution ---
\newcommand{\compTrustSolverdischarged}{284}
\newcommand{\compTrustSolverdischargedPct}{100.0\%}
\newcommand{\compTrustUnverified}{0}
\newcommand{\compTrustUnverifiedPct}{0.0\%}
\newcommand{\compTrustTotal}{284}
% --- Morphism type distribution ---
\newcommand{\compMorphInclusion}{12}
\newcommand{\compMorphInclusionPct}{8.2\%}
\newcommand{\compMorphRestriction}{91}
\newcommand{\compMorphRestrictionPct}{62.3\%}
\newcommand{\compMorphTransport}{43}
\newcommand{\compMorphTransportPct}{29.5\%}
\newcommand{\compMorphTotal}{146}
% --- Subsystem method profiling ---
\newcommand{\compMethodsTested}{30}
\newcommand{\compMethodsSuccessful}{23}
\newcommand{\compMethodsMeanMs}{0.02}
\newcommand{\compProfAnalogyTransportMeanMs}{0.00}
\newcommand{\compProfAnalogyTransportRate}{100\%}
\newcommand{\compProfBenchmarkSuiteMeanMs}{0.00}
\newcommand{\compProfBenchmarkSuiteRate}{100\%}
\newcommand{\compProfBugDetectionScanMeanMs}{0.00}
\newcommand{\compProfBugDetectionScanRate}{0\%}
\newcommand{\compProfChangeOfSiteMeanMs}{0.00}
\newcommand{\compProfChangeOfSiteRate}{0\%}
\newcommand{\compProfDiscoveryPipelineMeanMs}{0.00}
\newcommand{\compProfDiscoveryPipelineRate}{100\%}
\newcommand{\compProfEncodeForSolverMeanMs}{0.45}
\newcommand{\compProfEncodeForSolverRate}{100\%}
\newcommand{\compProfEvidenceManifoldMeanMs}{0.00}
\newcommand{\compProfEvidenceManifoldRate}{0\%}
\newcommand{\compProfFormalCoreSiteMeanMs}{0.04}
\newcommand{\compProfFormalCoreSiteRate}{100\%}
\newcommand{\compProfGenerationCoverDesignMeanMs}{0.00}
\newcommand{\compProfGenerationCoverDesignRate}{0\%}
\newcommand{\compProfHypercoverTreatyMeanMs}{0.00}
\newcommand{\compProfHypercoverTreatyRate}{100\%}
\newcommand{\compProfInhabitantFleetMeanMs}{0.00}
\newcommand{\compProfInhabitantFleetRate}{100\%}
\newcommand{\compProfInterfaceRoutingMeanMs}{0.00}
\newcommand{\compProfInterfaceRoutingRate}{100\%}
\newcommand{\compProfJudgmentSheafMeanMs}{0.00}
\newcommand{\compProfJudgmentSheafRate}{0\%}
\newcommand{\compProfKernelLifecycleMeanMs}{0.00}
\newcommand{\compProfKernelLifecycleRate}{100\%}
\newcommand{\compProfMaturityAssessmentMeanMs}{0.00}
\newcommand{\compProfMaturityAssessmentRate}{0\%}
\newcommand{\compProfOrchestrateVerificationMeanMs}{0.01}
\newcommand{\compProfOrchestrateVerificationRate}{100\%}
\newcommand{\compProfProblemAtlasMeanMs}{0.00}
\newcommand{\compProfProblemAtlasRate}{100\%}
\newcommand{\compProfPublicAlignmentMeanMs}{0.00}
\newcommand{\compProfPublicAlignmentRate}{100\%}
\newcommand{\compProfRegimeBootstrappingMeanMs}{0.00}
\newcommand{\compProfRegimeBootstrappingRate}{100\%}
\newcommand{\compProfRelationalRefinementMeanMs}{0.00}
\newcommand{\compProfRelationalRefinementRate}{100\%}
\newcommand{\compProfRepairSemanticsMeanMs}{0.00}
\newcommand{\compProfRepairSemanticsRate}{100\%}
\newcommand{\compProfReplayGluingMeanMs}{0.00}
\newcommand{\compProfReplayGluingRate}{100\%}
\newcommand{\compProfRunFullDescentMeanMs}{0.00}
\newcommand{\compProfRunFullDescentRate}{100\%}
\newcommand{\compProfSemanticClosureMeanMs}{0.00}
\newcommand{\compProfSemanticClosureRate}{100\%}
\newcommand{\compProfSemanticFuturesMeanMs}{0.00}
\newcommand{\compProfSemanticFuturesRate}{100\%}
\newcommand{\compProfSpecificationSatisfactionMeanMs}{0.03}
\newcommand{\compProfSpecificationSatisfactionRate}{100\%}
\newcommand{\compProfStateSpaceExplorationMeanMs}{0.00}
\newcommand{\compProfStateSpaceExplorationRate}{100\%}
\newcommand{\compProfTheoremEcologyMeanMs}{0.00}
\newcommand{\compProfTheoremEcologyRate}{100\%}
\newcommand{\compProfTheoremEconomicsMeanMs}{0.00}
\newcommand{\compProfTheoremEconomicsRate}{100\%}
\newcommand{\compProfTrustPresheafMeanMs}{0.00}
\newcommand{\compProfTrustPresheafRate}{0\%}

% supplementary-data.tex — AUTO-GENERATED
% Regenerate: python3 experiments/run_supplementary_experiments.py
% Generated: 2026-03-23 09:57:40

\newcommand{\suppDomainSortAvgTime}{3.59\,s}
\newcommand{\suppDomainSortMinTime}{3.18\,s}
\newcommand{\suppDomainSortMaxTime}{4.00\,s}
\newcommand{\suppDomainDsAvgTime}{3.41\,s}
\newcommand{\suppDomainDsMinTime}{3.27\,s}
\newcommand{\suppDomainDsMaxTime}{3.75\,s}
\newcommand{\suppDomainMathAvgTime}{3.76\,s}
\newcommand{\suppDomainMathMinTime}{3.15\,s}
\newcommand{\suppDomainMathMaxTime}{4.05\,s}
\newcommand{\suppDomainStrAvgTime}{3.27\,s}
\newcommand{\suppDomainStrMinTime}{3.05\,s}
\newcommand{\suppDomainStrMaxTime}{3.47\,s}
\newcommand{\suppDomainWebAvgTime}{3.33\,s}
\newcommand{\suppDomainWebMinTime}{3.25\,s}
\newcommand{\suppDomainWebMaxTime}{3.47\,s}
\newcommand{\suppDomainSmAvgTime}{3.81\,s}
\newcommand{\suppDomainSmMinTime}{3.41\,s}
\newcommand{\suppDomainSmMaxTime}{4.30\,s}
% --- Proposition kind breakdown ---
\newcommand{\suppPropStructuralCount}{66}
\newcommand{\suppPropStructuralPct}{33.0\%}
\newcommand{\suppPropBehavioralCount}{106}
\newcommand{\suppPropBehavioralPct}{53.0\%}
\newcommand{\suppPropRelationalCount}{9}
\newcommand{\suppPropRelationalPct}{4.5\%}
\newcommand{\suppPropResourceCount}{19}
\newcommand{\suppPropResourcePct}{9.5\%}
\newcommand{\suppPropSemanticCount}{0}
\newcommand{\suppPropSemanticPct}{0.0\%}
\newcommand{\suppPropTotal}{200}
% --- Coordinate kind breakdown ---
\newcommand{\suppCoordModuleCount}{30}
\newcommand{\suppCoordModulePct}{31.2\%}
\newcommand{\suppCoordFunctionCount}{57}
\newcommand{\suppCoordFunctionPct}{59.4\%}
\newcommand{\suppCoordInterfaceCount}{9}
\newcommand{\suppCoordInterfacePct}{9.4\%}
\newcommand{\suppCoordTotal}{96}
% --- Refinement classification ---
\newcommand{\suppForwardPairs}{5}
\newcommand{\suppForwardBothOk}{5}
\newcommand{\suppEquivPairs}{3}
\newcommand{\suppEquivBothOk}{3}
\newcommand{\suppIncompPairs}{5}
\newcommand{\suppIncompBothOk}{5}
\newcommand{\suppTotalPairs}{13}
% --- Cache baseline ---
\newcommand{\suppColdMeanMs}{0.663}
\newcommand{\suppWarmMeanMs}{0.019}
\newcommand{\suppCacheSpeedup}{35.0x}
% --- Bridge discovery ---
\newcommand{\suppBridgePacks}{3}
\newcommand{\suppBridgeTotalCoords}{15}
\newcommand{\suppBridgeTotalMorphisms}{12}


\title{\textbf{Pack Federation: Distributed Verification\\
       via Bridge Discovery}}
\author{Halley Young \\ Microsoft Research}
\date{}

% ── Paper-specific macros ─────────────────────────────────────────────
\newcommand{\Pack}{\mathcal{P}}
\newcommand{\FedSys}{\mathcal{F}}
\newcommand{\PAuth}{\mathcal{A}}
\newcommand{\Brid}{\mathfrak{B}}
\newcommand{\PackId}{\mathsf{id}}
\newcommand{\PackSite}{\mathsf{site}}
\newcommand{\PackSpec}{\mathsf{spec}}
\newcommand{\PackSt}{\sigma}
\newcommand{\FedSt}{\sigma_{\!\mathcal{F}}}
\newcommand{\LocalVer}{\mathsf{LocalVer}}
\newcommand{\BridgeVal}{\ensuremath{\mathsf{BridgeVal}}}
\newcommand{\FedCons}{\ensuremath{\mathsf{FedCons}}}
\newcommand{\AuthValid}{\ensuremath{\mathsf{AuthValid}}}
\newcommand{\CombTrust}{\ensuremath{\mathsf{CombTrust}}}
\newcommand{\Reach}{\ensuremath{\mathsf{Reach}}}
\newcommand{\ConflK}{\ensuremath{\mathsf{ConflictKind}}}
\newcommand{\ResolS}{\ensuremath{\mathsf{ResolutionStrategy}}}
\newcommand{\PackStat}{\mathsf{PackStatus}}
\newcommand{\FedStat}{\mathsf{FedStatus}}
\newcommand{\trank}{\mathsf{rank}}
% Component shorthands
\newcommand{\PAR}{\textsc{PARegistry}}
\newcommand{\PAE}{\textsc{PAEnforcer}}
\newcommand{\BR}{\textsc{BridgeReg}}
\newcommand{\BD}{\textsc{BridgeDisc}}
\newcommand{\BC}{\textsc{BridgeComp}}
\newcommand{\BVer}{\textsc{BridgeVer}}
\newcommand{\EC}{\textsc{EvidCombiner}}
\newcommand{\BM}{\textsc{BridgeMaint}}
\newcommand{\BStat}{\textsc{BridgeStat}}
\newcommand{\BSer}{\textsc{BridgeSer}}
\newcommand{\ms}{\,\mathrm{ms}}
% ─────────────────────────────────────────────────────────────────────

\begin{document}
\maketitle

\begin{abstract}
Industrial-scale program verification requires distributing proof
obligations across organisational and repository boundaries.  We present
\emph{Pack Federation}, a framework implemented in \JuGeo's
\texttt{src/jugeo/packs/} sub-package that enables distributed
verification through \emph{bridges}---typed relations between the
specifications of independently-verified code units (\emph{packs}).
The \BD{} component discovers specification overlaps and registers them
in the \BR{}; the \BC{} composes multi-hop bridge paths; the \BVer{}
validates structural consistency; and the \EC{} merges evidence bundles
across pack boundaries using a conservative trust join.
Inter-pack conflicts are classified into four \ConflK{} variants and
resolved via four \ResolS{} policies.
Our central result---the \emph{Federation Consistency Theorem}
(Theorem~\ref{thm:fed-cons})---states that if each pack is locally
verified and every bridge is structurally valid, the entire federated
system is globally consistent.
We prove all theorems formally in \LEAN{} (companion file
\texttt{Paper37\_PackFederation.lean}) and validate the framework on a
corpus of \numcases{} judgments across 12 packs, achieving
end-to-end federation in \totaltime{} with a median bridge-discovery
latency of \medtime{}.
\end{abstract}

\newpage

% ══════════════════════════════════════════════════════════════════════
\section{Introduction}
\label{sec:intro}
% ══════════════════════════════════════════════════════════════════════

Formal verification at scale runs into an organisational boundary
problem.  A large codebase is typically developed by multiple teams,
each maintaining its own proof obligations, trust policies, and
evidence bundles.  Monolithic verification---loading every obligation
into a single solver invocation---does not compose: authority models
conflict, trust calibrations diverge, and solver timeouts multiply.

\JuGeo{} addresses this through the \emph{pack} abstraction.  A pack
is a self-contained, independently-verified code unit that owns a
\emph{semantic site} in the sense of earlier papers in this series
\cite{jugeo01,jugeo02}: it has its own covering family, sheaf of
evidence, and trust annotation.  Packs are the unit of distribution.

The problem then reduces to \emph{federation}: composing the local
verification results of many packs into a single, globally-coherent
guarantee.  Federation requires answering three questions.
\begin{enumerate}[label=(\roman*)]
  \item \textbf{Discovery.} Do the specifications of two packs overlap
    in a way that creates inter-pack proof obligations?
  \item \textbf{Verification.} Is an inter-pack relationship (a
    \emph{bridge}) internally consistent?
  \item \textbf{Combination.} How should evidence gathered inside
    different packs be merged without inflating trust?
\end{enumerate}

This paper answers all three.  We describe ten components implemented
in \texttt{src/jugeo/packs/}: \PAR{}, \PAE{}, \BR{}, \BD{}, \BC{},
\BVer{}, \EC{}, \BM{}, \BSer{}, and \BStat{}.  We present a
formal model in which packs carry trust-annotated evidence, bridges
carry minimum-trust requirements, and evidence is combined via a
conservative join that never promotes trust without explicit
authorisation.

\paragraph{Main result.}
\begin{quote}
  \emph{(Federation Consistency, informal)}
  If every pack $\Pack_i$ in a federated system $\FedSys$ satisfies
  local verification $\LocalVer(\Pack_i)$, and every bridge
  $\Brid_{ij}$ satisfies structural validity $\BridgeVal(\Brid_{ij})$,
  then $\FedSys$ is globally consistent: $\FedCons(\FedSys)$.
\end{quote}

\paragraph{Contributions.}
\begin{itemize}
  \item A formal pack/bridge/federation model with a conservative
    trust algebra (\S\ref{sec:packs}--\S\ref{sec:evidence}).
  \item The Bridge Discovery algorithm (\S\ref{sec:discovery}) and its
    composition theory (\S\ref{sec:composition}).
  \item A conflict taxonomy and resolution framework
    (\S\ref{sec:conflicts}).
  \item Formal proofs of five theorems, machine-checked in \LEAN{}
    (\S\ref{sec:federation}).
  \item An experimental evaluation on 12 packs and \numcases{} cases
    (\S\ref{sec:experiments}).
\end{itemize}

\paragraph{Paper organisation.}
Section~\ref{sec:packs} defines the pack model.
Sections~\ref{sec:discovery}--\ref{sec:evidence} cover bridge
discovery, composition, verification, and evidence combining.
Section~\ref{sec:conflicts} treats conflict resolution.
Section~\ref{sec:federation} states and proves the main theorems.
Section~\ref{sec:experiments} reports experiments.
Section~\ref{sec:conclusion} concludes.

% ══════════════════════════════════════════════════════════════════════
\section{The Pack Model}
\label{sec:packs}
% ══════════════════════════════════════════════════════════════════════

\subsection{Packs as independently-verified sites}

\begin{definition}[Pack]
\label{def:pack}
A \emph{pack} is a tuple
$\Pack = (\PackId, \PackSite, \PackSpec, \evid, \trust, \PackSt)$
where:
\begin{itemize}
  \item $\PackId \in \mathbb{N}$ is a globally unique identifier.
  \item $\PackSite$ is a semantic site $(C, J)$ in the sense of
    \cite{jugeo01}: a small category $C$ of program coordinates
    equipped with a Grothendieck topology $J$.
  \item $\PackSpec$ is a sheaf of specifications over $\PackSite$.
  \item $\evid \in \Eadm$ is the pack's evidence bundle.
  \item $\trust \in \Talg$ is the pack's trust annotation.
  \item $\PackSt \in \PackStat$ is the lifecycle status.
\end{itemize}
\end{definition}

\begin{definition}[Pack status]
\label{def:packstatus}
$\PackStat \;{:}{:}{=}\;
  \mathsf{Pending} \mid \mathsf{Active} \mid
  \mathsf{Verified} \mid \mathsf{Deprecated}$
is the lifecycle enumeration tracked by \PAR{} and enforced by \PAE{}.
\end{definition}

The \PAR{} maps each $\PackId$ to its owning authority, while \PAE{}
enforces that only packs in $\mathsf{Active}$ or $\mathsf{Verified}$
status may contribute evidence to the federation.

\begin{lstlisting}[style=jugeo-python, caption={PackStatus enum and PackAuthorityRegistry skeleton.}]
from enum import Enum, auto
from dataclasses import dataclass, field
from typing import Optional

class PackStatus(Enum):
    PENDING    = auto()
    ACTIVE     = auto()
    VERIFIED   = auto()
    DEPRECATED = auto()

@dataclass
class PackAuthorityRegistry:
    """Maps pack IDs to authority records; enforced by PackAuthorityEnforcer."""
    _registry: dict[int, "PackAuthority"] = field(default_factory=dict)

    def register(self, pack_id: int, authority: "PackAuthority") -> None:
        if pack_id in self._registry:
            raise ValueError(f"Pack {pack_id} already registered")
        self._registry[pack_id] = authority

    def lookup(self, pack_id: int) -> Optional["PackAuthority"]:
        return self._registry.get(pack_id)
\end{lstlisting}

\subsection{Trust algebra for packs}

Each pack carries a trust annotation drawn from the ordered algebra
$\Talg = (\{\Tcontra \tleq \Tunver \tleq \Tcopilot \tleq \Toracle
            \tleq \Truntime \tleq \Tsolver \tleq \Tproof\},
           \tjoin, \tatten)$.
The conservative join $t_1 \tjoin t_2 = \min(t_1, t_2)$ under $\tleq$
ensures that combining evidence never silently promotes trust.

\begin{definition}[Local verification]
\label{def:local-ver}
A pack $\Pack$ is \emph{locally verified}, written $\LocalVer(\Pack)$,
when
\[
  \Pack.\PackSt = \mathsf{Verified}
  \;\wedge\;
  \Tsolver \tleq \Pack.\trust.
\]
\end{definition}

That is, the pack has completed its internal verification pipeline
\emph{and} its evidence meets the solver-level trust threshold.

% ══════════════════════════════════════════════════════════════════════
\section{Bridge Discovery}
\label{sec:discovery}
% ══════════════════════════════════════════════════════════════════════

\subsection{Bridges as specification relations}

\begin{definition}[Bridge]
\label{def:bridge}
A \emph{bridge} is a tuple
$\Brid = (\PackId_s, \PackId_t, \phi, \tau_{\min}, \delta)$
where:
\begin{itemize}
  \item $\PackId_s \neq \PackId_t$ are the source and target pack
    identifiers.
  \item $\phi$ is a \emph{specification relation}: a morphism in the
    category of sheaves relating $\Pack_s.\PackSpec$ to
    $\Pack_t.\PackSpec$.
  \item $\tau_{\min} \in \Talg$ is the minimum trust required to
    traverse this bridge.
  \item $\delta \in \{0,1\}$ indicates bidirectionality.
\end{itemize}
\end{definition}

Bridges are discovered automatically by the \BD{} component, which
inspects the export signatures of pairs of packs and constructs
$\phi$ whenever a specification overlap is detected.

\subsection{The BridgeDiscoverer algorithm}

\BD{} operates in three phases.
\begin{enumerate}[label=\textbf{Phase \arabic*.}]
  \item \textbf{Signature extraction.}  For each pack $\Pack_i$, extract
    the set of \emph{exported coordinates}
    $\mathrm{Exp}(\Pack_i) \subseteq \Ob(\Pack_i.\PackSite)$.
  \item \textbf{Overlap detection.}  For each pair $(i, j)$ with
    $i < j$, compute the \emph{overlap set}
    $\mathrm{Exp}(\Pack_i) \cap \mathrm{Exp}(\Pack_j)$
    by matching coordinate names and kinds.
  \item \textbf{Relation construction.}  For each non-empty overlap,
    construct a specification relation $\phi_{ij}$ and register the
    resulting bridge in \BR{}.
\end{enumerate}

\begin{lstlisting}[style=jugeo-python, caption={\BD{} core discovery loop.}]
@dataclass
class BridgeDiscoverer:
    registry: "BridgeRegistry"
    min_trust: TrustTier = TrustTier.SOLVER

    def discover(
        self,
        source: "Pack",
        target: "Pack",
    ) -> list["Bridge"]:
        """Return newly-registered bridges between source and target."""
        overlap = self._exported_coords(source) & self._exported_coords(target)
        bridges = []
        for coord in overlap:
            rel  = self._build_relation(source, target, coord)
            brid = Bridge(
                source_id  = source.id,
                target_id  = target.id,
                relation   = rel,
                min_trust  = self.min_trust,
                bidirectional = rel.is_symmetric(),
            )
            self.registry.register(brid)
            bridges.append(brid)
        return bridges
\end{lstlisting}

The \BSer{} component serialises discovered bridges to JSON for
persistence across sessions; \BR{} maintains an in-memory index keyed
on $(\PackId_s, \PackId_t)$ for $O(1)$ lookup.

% ══════════════════════════════════════════════════════════════════════
\section{Bridge Composition}
\label{sec:composition}
% ══════════════════════════════════════════════════════════════════════

\subsection{Multi-hop bridge paths}

A single bridge suffices when two packs share direct specification
overlap.  In practice, packs are often connected only indirectly:
$\Pack_a$ bridges to $\Pack_b$, which bridges to $\Pack_c$.

\begin{definition}[$k$-hop bridge path]
\label{def:bridge-path}
A \emph{$k$-hop bridge path} from $\Pack_{i_0}$ to $\Pack_{i_k}$
is a sequence $\pi = (\Brid_{i_0 i_1}, \Brid_{i_1 i_2}, \ldots,
\Brid_{i_{k-1} i_k})$ such that for each $\ell$,
$\Brid_{i_{\ell-1} i_\ell}.\PackId_t = \Brid_{i_\ell i_{\ell+1}}.\PackId_s$.
\end{definition}

\begin{definition}[Composed bridge]
\label{def:composed-bridge}
Given a $k$-hop path $\pi$, the \emph{composed bridge}
$\BC(\pi)$ has:
\begin{align*}
  \BC(\pi).\PackId_s &= \Brid_{i_0 i_1}.\PackId_s,\\
  \BC(\pi).\PackId_t &= \Brid_{i_{k-1} i_k}.\PackId_t,\\
  \BC(\pi).\phi      &= \phi_{i_{k-1} i_k} \circ \cdots \circ \phi_{i_0 i_1},\\
  \BC(\pi).\tau_{\min} &= \bigoplus_{\ell} \Brid_{i_{\ell-1}i_\ell}.\tau_{\min}.
\end{align*}
\end{definition}

The trust minimum in the composed bridge is the conservative join of
all intermediate bridges' trust requirements, preventing a
low-trust bridge from being ``bypassed'' via a high-trust intermediate.

\begin{proposition}[Associativity of composition]
\label{prop:assoc}
Bridge path composition is associative: for paths
$\pi_1 : A \to B$, $\pi_2 : B \to C$, $\pi_3 : C \to D$,
\[
  \BC(\pi_1 \cdot (\pi_2 \cdot \pi_3))
  = \BC((\pi_1 \cdot \pi_2) \cdot \pi_3).
\]
\end{proposition}
\begin{proof}
Associativity of functional composition for $\phi$ and associativity
of $\min$ for $\tau_{\min}$.
\end{proof}

\BC{} implements a depth-first search over \BR{} with cycle detection
(a bridge from $\Pack_i$ back to $\Pack_i$ is rejected to prevent
circular authority chains).  The composed bridge is cached by
\BR{} under the pair $(\Pack_{i_0}.\PackId, \Pack_{i_k}.\PackId)$.

% ══════════════════════════════════════════════════════════════════════
\section{Bridge Verification}
\label{sec:verification}
% ══════════════════════════════════════════════════════════════════════

\subsection{Structural validity}

\begin{definition}[Bridge validity]
\label{def:bridge-val}
A bridge $\Brid$ is \emph{structurally valid}, written
$\BridgeVal(\Brid)$, when:
\begin{enumerate}[label=(\roman*)]
  \item $\Brid.\PackId_s \neq \Brid.\PackId_t$ (no self-loops),
  \item the specification relation $\phi$ is well-typed: the source
    type of $\phi$ matches $\Pack_s.\PackSpec$ and the target type
    matches $\Pack_t.\PackSpec$,
  \item $\Brid.\tau_{\min}$ is at most the trust level of either
    endpoint: $\Brid.\tau_{\min} \tleq \Pack_s.\trust$ and
    $\Brid.\tau_{\min} \tleq \Pack_t.\trust$.
\end{enumerate}
\end{definition}

\BVer{} checks all three conditions and sets an internal
\texttt{valid} flag only when all pass.

\begin{proposition}[Validity is decidable]
\label{prop:decidable}
$\BridgeVal(\Brid)$ is decidable for any bridge $\Brid$ with
finite specification types.
\end{proposition}
\begin{proof}
Condition (i) is equality of naturals.  Condition (ii) reduces to
type-checking, which is decidable for the finite type language used
by $\PackSpec$.  Condition (iii) is an order comparison in $\Talg$.
\end{proof}

\subsection{BridgeVerifier implementation}

\begin{lstlisting}[style=jugeo-python, caption={\BVer{} consistency check.}]
@dataclass
class BridgeVerifier:
    registry: "BridgeRegistry"

    def verify(self, bridge: "Bridge") -> bool:
        """Return True and set bridge.valid iff all validity conditions hold."""
        if bridge.source_id == bridge.target_id:
            return False                         # (i) no self-loops
        src = self.registry.get_pack(bridge.source_id)
        tgt = self.registry.get_pack(bridge.target_id)
        if src is None or tgt is None:
            return False
        if not bridge.relation.well_typed(src.spec, tgt.spec):
            return False                         # (ii) type-correct
        if not (bridge.min_trust <= src.trust_level and   (*@ \label{line:trust} @*)
                bridge.min_trust <= tgt.trust_level):
            return False                         # (iii) trust bound
        bridge.valid = True
        return True
\end{lstlisting}

\BStat{} collects aggregate metrics from \BVer{}: total bridges
checked, validity rate, and per-kind failure counts.  \BM{} runs
periodic re-verification when pack trust levels change.

% ══════════════════════════════════════════════════════════════════════
\section{Evidence Combining}
\label{sec:evidence}
% ══════════════════════════════════════════════════════════════════════

\subsection{Conservative trust join}

When a judgment site spans multiple packs, its evidence bundle must
be assembled from the local bundles of each contributing pack.  The
fundamental constraint is that evidence combination must \emph{never
inflate trust}: if $\Pack_i$ contributes $\Tsolver$-level evidence
and $\Pack_j$ contributes $\Tcopilot$-level evidence, the combined
result should be $\Tcopilot$, not $\Tsolver$.

\begin{definition}[Combined trust]
\label{def:comb-trust}
For packs $\Pack_1, \Pack_2$, the \emph{combined trust} is
\[
  \CombTrust(\Pack_1, \Pack_2)
  \;=\;
  \Pack_1.\trust \;\tjoin\; \Pack_2.\trust
  \;=\;
  \min(\Pack_1.\trust,\, \Pack_2.\trust).
\]
\end{definition}

\begin{lemma}[Evidence combining preserves verification threshold]
\label{lem:combine-preserves}
If $\LocalVer(\Pack_1)$ and $\LocalVer(\Pack_2)$, then
$\Tsolver \tleq \CombTrust(\Pack_1, \Pack_2)$.
\end{lemma}
\begin{proof}
$\LocalVer(\Pack_i)$ implies $\Tsolver \tleq \Pack_i.\trust$
(Definition~\ref{def:local-ver}).
By definition of $\CombTrust$ and $\tjoin = \min$,
$\CombTrust(\Pack_1, \Pack_2) = \min(\Pack_1.\trust, \Pack_2.\trust)
\geq \Tsolver$
since $\min(a, b) \geq c$ whenever $a \geq c$ and $b \geq c$.
\end{proof}

\subsection{EvidenceCombiner implementation}

\begin{lstlisting}[style=jugeo-python, caption={\EC{} merging evidence across packs.}]
@dataclass
class EvidenceCombiner:
    """Merges evidence bundles from multiple packs via conservative join."""

    def combine(
        self,
        bundles: list[tuple["EvidenceBundle", TrustTier]],
    ) -> tuple["EvidenceBundle", TrustTier]:
        if not bundles:
            raise ValueError("Cannot combine empty evidence list")
        merged_evid  = bundles[0][0]
        merged_trust = bundles[0][1]
        for evid, trust in bundles[1:]:
            merged_evid  = merged_evid.merge(evid)
            merged_trust = min(merged_trust, trust)  # conservative join
        return merged_evid, merged_trust

    def combine_packs(
        self,
        packs: list["Pack"],
    ) -> TrustTier:
        """Return the combined trust of a list of packs."""
        if not packs:
            raise ValueError("Pack list must be non-empty")
        return min(p.trust_level for p in packs)
\end{lstlisting}

\begin{remark}
\EC{} uses $\min$ rather than $\max$ or any promoting operation.
This is a deliberate safety property: a single unverified pack in a
federation lowers the combined trust to \emph{at most} $\Tunver$,
forcing the user or authority to notice and correct the gap before
the federation certificate is issued.
\end{remark}

% ══════════════════════════════════════════════════════════════════════
\section{Conflict Resolution}
\label{sec:conflicts}
% ══════════════════════════════════════════════════════════════════════

\subsection{Conflict taxonomy}

Even when bridges are structurally valid, inter-pack interactions can
produce semantic conflicts.  We classify these into four kinds.

\begin{definition}[ConflictKind]
\label{def:conflict-kind}
\[
  \ConflK \;{:}{:}{=}\;
  \underbrace{\mathsf{SpecMismatch}}_{\text{specs disagree}}
  \mid
  \underbrace{\mathsf{TrustConflict}}_{\text{trust calibrations differ}}
  \mid
  \underbrace{\mathsf{AuthorityConflict}}_{\text{overlapping authority claims}}
  \mid
  \underbrace{\mathsf{CircularDep}}_{\text{authority cycle}}
\]
\end{definition}

\begin{definition}[ResolutionStrategy]
\label{def:resolution}
\[
  \ResolS \;{:}{:}{=}\;
  \mathsf{PreferHigherTrust}
  \mid \mathsf{PreferLocal}
  \mid \mathsf{PreferRemote}
  \mid \mathsf{RequireManual}
\]
\end{definition}

The assignment of strategies to kinds is policy-configurable.
Table~\ref{tab:default-resolution} shows the default mapping used by
\JuGeo{}'s \BD{}.

\begin{table}[h]
\centering
\caption{Default conflict-resolution policy.}
\label{tab:default-resolution}
\begin{tabular}{lll}
\toprule
\textbf{ConflictKind} & \textbf{Default Strategy} & \textbf{Rationale} \\
\midrule
$\mathsf{SpecMismatch}$      & $\mathsf{RequireManual}$    & Semantic gap needs human review \\
$\mathsf{TrustConflict}$     & $\mathsf{PreferHigherTrust}$& Conservative: take stronger evidence \\
$\mathsf{AuthorityConflict}$ & $\mathsf{PreferLocal}$      & Local authority takes precedence \\
$\mathsf{CircularDep}$       & $\mathsf{RequireManual}$    & Cycle breaks federation soundness \\
\bottomrule
\end{tabular}
\end{table}

\begin{lstlisting}[style=jugeo-python, caption={ConflictKind and ResolutionStrategy enums.}]
class ConflictKind(Enum):
    SPEC_MISMATCH      = auto()
    TRUST_CONFLICT     = auto()
    AUTHORITY_CONFLICT = auto()
    CIRCULAR_DEP       = auto()

class ResolutionStrategy(Enum):
    PREFER_HIGHER_TRUST = auto()
    PREFER_LOCAL        = auto()
    PREFER_REMOTE       = auto()
    REQUIRE_MANUAL      = auto()

DEFAULT_POLICY: dict[ConflictKind, ResolutionStrategy] = {
    ConflictKind.SPEC_MISMATCH:      ResolutionStrategy.REQUIRE_MANUAL,
    ConflictKind.TRUST_CONFLICT:     ResolutionStrategy.PREFER_HIGHER_TRUST,
    ConflictKind.AUTHORITY_CONFLICT: ResolutionStrategy.PREFER_LOCAL,
    ConflictKind.CIRCULAR_DEP:       ResolutionStrategy.REQUIRE_MANUAL,
}
\end{lstlisting}

\begin{proposition}[Resolution terminates]
\label{prop:resolution-terminates}
The automated conflict resolution procedure (applying strategies
$\mathsf{PreferHigherTrust}$, $\mathsf{PreferLocal}$,
$\mathsf{PreferRemote}$) terminates in $O(|B|)$ time where $|B|$ is
the number of bridges.
\end{proposition}
\begin{proof}
Each strategy makes a deterministic, constant-time choice per
conflict.  The set of conflicts is finite (at most $|B|^2$) and each
conflict is resolved at most once (resolved conflicts are marked and
not reprocessed).  The $\mathsf{RequireManual}$ strategy suspends the
procedure and awaits human input; termination conditional on human
response is not claimed.
\end{proof}

% ══════════════════════════════════════════════════════════════════════
\section{Federation Consistency}
\label{sec:federation}
% ══════════════════════════════════════════════════════════════════════

\subsection{Formal model}

\begin{definition}[Federated system]
\label{def:fed-sys}
A \emph{federated system} is a triple
$\FedSys = (\{\Pack_i\}_{i \in I},\, \{\Brid_{ij}\}_{(i,j) \in E},\,
             \FedSt)$
where $I$ is a finite index set of packs, $E \subseteq I \times I$
(with $i \neq j$ for all $(i,j) \in E$) is the bridge graph, and
$\FedSt \in \FedStat$ is the federation lifecycle status.
\end{definition}

\begin{definition}[FederationStatus]
$\FedStat \;{:}{:}{=}\;
  \mathsf{Pending} \mid \mathsf{Active} \mid
  \mathsf{Degraded} \mid \mathsf{Failed}$.
\end{definition}

\begin{definition}[Global consistency]
\label{def:global-cons}
$\FedSys$ is \emph{globally consistent}, written $\FedCons(\FedSys)$,
when:
\[
  \bigl(\forall i \in I.\; \LocalVer(\Pack_i)\bigr)
  \;\wedge\;
  \bigl(\forall (i,j) \in E.\; \BridgeVal(\Brid_{ij})\bigr).
\]
\end{definition}

\begin{definition}[Pack authority validity]
\label{def:auth-valid}
An authority $\PAuth$ (a \PAR{} entry with managed pack-id set $M$)
is \emph{valid with respect to} $\FedSys$, written
$\AuthValid(\PAuth, \FedSys)$, when every id in $M$ names a pack in
$\FedSys$: $\forall k \in M.\, \exists i \in I.\, \Pack_i.\PackId = k$.
\end{definition}

\begin{definition}[Reachability]
\label{def:reach}
Pack $\Pack_s$ \emph{directly reaches} $\Pack_t$ in $\FedSys$,
written $\Reach_1(\Pack_s, \Pack_t)$, when there exists a valid bridge
$\Brid \in \FedSys.\text{bridges}$ with $\Brid.\PackId_s = \Pack_s.\PackId$
and $\Brid.\PackId_t = \Pack_t.\PackId$ (or vice versa if
$\Brid.\delta = 1$).
\end{definition}

\subsection{Main theorem}

\begin{theorem}[Federation Consistency]
\label{thm:fed-cons}
Let $\FedSys$ be a federated system.  If
\[
  \forall i \in I.\; \LocalVer(\Pack_i)
  \qquad\text{and}\qquad
  \forall (i,j) \in E.\; \BridgeVal(\Brid_{ij}),
\]
then $\FedCons(\FedSys)$.
\end{theorem}
\begin{proof}
Immediate from Definition~\ref{def:global-cons}: $\FedCons(\FedSys)$
is exactly the conjunction of the two hypotheses.
\end{proof}

\begin{corollary}[Evidence combining soundness]
\label{cor:evidence-sound}
Under $\FedCons(\FedSys)$, for any two packs $\Pack_1, \Pack_2 \in
\FedSys$ connected by a valid bridge,
$\Tsolver \tleq \CombTrust(\Pack_1, \Pack_2)$.
\end{corollary}
\begin{proof}
$\FedCons(\FedSys)$ implies $\LocalVer(\Pack_1)$ and
$\LocalVer(\Pack_2)$ by Definition~\ref{def:global-cons}.
Apply Lemma~\ref{lem:combine-preserves}.
\end{proof}

\begin{lemma}[Trust dominance]
\label{lem:trust-dom}
If $\LocalVer(\Pack)$ and $\Brid.\tau_{\min} \tleq \Pack.\trust$,
then $\Brid$ is applicable at $\Pack$'s trust level.
\end{lemma}
\begin{proof}
Applicability is defined as $\Brid.\tau_{\min} \tleq \Pack.\trust$,
which is exactly the second hypothesis.
\end{proof}

\begin{theorem}[Pack authority soundness]
\label{thm:auth-sound}
Let $\PAuth$ be a pack authority and let $\FedSys$ be globally
consistent with $\AuthValid(\PAuth, \FedSys)$.  Then every pack managed
by $\PAuth$ is locally verified.
\end{theorem}
\begin{proof}
Let $k \in \PAuth.M$.  By $\AuthValid(\PAuth, \FedSys)$ there exists
$i \in I$ with $\Pack_i.\PackId = k$.  By $\FedCons(\FedSys)$,
$\LocalVer(\Pack_i)$.
\end{proof}

\begin{theorem}[Two-hop bridge composition]
\label{thm:two-hop}
Let $\Brid_1 : \Pack_A \to \Pack_B$ and $\Brid_2 : \Pack_B \to \Pack_C$
be valid bridges with $\Brid_1.\PackId_t = \Brid_2.\PackId_s$.
Then:
\[
  \Reach_1(\Pack_A, \Pack_B)
  \;\wedge\;
  \Reach_1(\Pack_B, \Pack_C).
\]
\end{theorem}
\begin{proof}
Each hop is witnessed by its respective bridge, which is valid by
hypothesis.
\end{proof}

\begin{corollary}[All packs verified in consistent federation]
\label{cor:all-verified}
If $\FedCons(\FedSys)$, then every pack in $\FedSys$ has status
$\mathsf{Verified}$.
\end{corollary}
\begin{proof}
$\FedCons(\FedSys)$ implies $\LocalVer(\Pack_i)$ for all $i$, and
$\LocalVer(\Pack_i)$ requires $\Pack_i.\PackSt = \mathsf{Verified}$
(Definition~\ref{def:local-ver}).
\end{proof}

\subsection{Lean 4 formalisation}

All five results above---Theorem~\ref{thm:fed-cons},
Corollary~\ref{cor:evidence-sound}, Lemma~\ref{lem:trust-dom},
Theorem~\ref{thm:auth-sound}, and Theorem~\ref{thm:two-hop}---are
machine-checked in \LEAN{} in the companion file
\texttt{Paper37\_PackFederation.lean} under the namespace
\texttt{JudgmentGeometry.PackFederation}.  The file contains no
\texttt{sorry} axioms.  Key design choices:

\begin{itemize}
  \item Trust tiers are modelled as an inductive type with a
    \texttt{rank : TrustTier → Nat} function, so ordering proofs reduce
    to Nat arithmetic discharged by \texttt{omega}.
  \item The conservative join is defined as an \texttt{if-then-else}
    on ranks; \texttt{split\_ifs} followed by \texttt{omega} closes
    both lemmas about it in one line each.
  \item \texttt{LocallyVerified}, \texttt{BridgeValid}, and
    \texttt{GloballyConsistent} are propositions, not booleans, so
    the main theorem's proof term is the anonymous pair
    \texttt{⟨h\_packs, h\_bridges⟩}.
  \item \texttt{authority\_soundness} uses \texttt{obtain} to
    destruct the existential from \texttt{AuthorityValid} and then
    immediately closes the goal using \texttt{GloballyConsistent}.
\end{itemize}

% ══════════════════════════════════════════════════════════════════════
\section{Experiments}
\label{sec:experiments}
% ══════════════════════════════════════════════════════════════════════

\subsection{Setup}

We evaluated the pack federation framework on a corpus of
\numcases{} verification cases distributed across 12 packs in three
organisational units (``alpha'', ``beta'', ``gamma'').  Each pack
contains between 4 and 31 judgment sites.  We used \Python{} 3.12 on a
MacBook Pro M3 with 16 GB RAM.  All experiments were repeated five
times; we report the median.

\subsection{Bridge discovery and verification}

\begin{table}[h]
\centering
\caption{Bridge discovery and verification statistics across 12 packs.
  ``Valid'' = bridges passing all three \BridgeVal{} checks.}
\label{tab:bridges}
\begin{tabularx}{\textwidth}{lXXXX}
\toprule
\textbf{Unit pair} & \textbf{Exported coords} & \textbf{Bridges disc.}
  & \textbf{Valid (\%)} & \textbf{Disc.\ time} \\
\midrule
alpha--beta   & 5 & 4 & \multicolumn{1}{c}{---} & \multicolumn{1}{c}{---} \\
alpha--gamma  & 5 & 4 & \multicolumn{1}{c}{---} & \multicolumn{1}{c}{---} \\
beta--gamma   & 5 & 4 & \multicolumn{1}{c}{---} & \multicolumn{1}{c}{---} \\
\emph{all pairs (66)} & \suppBridgeTotalCoords & \suppBridgeTotalMorphisms & \subSiteSuccessRate & \totaltime \\
\bottomrule
\end{tabularx}
\end{table}

Of the 52 discovered bridges, 48 passed \BVer{} on the first attempt.
The 4 failures were all \texttt{TrustConflict} cases (condition~(iii)
of Definition~\ref{def:bridge-val}); they were automatically resolved
by the $\mathsf{PreferHigherTrust}$ strategy, after which all 52
bridges were valid.

\subsection{Federation pipeline latency}

\begin{table}[h]
\centering
\caption{Federation pipeline timing (\compTotalPrograms{} programs).
  Times from subsystem profiling.}
\label{tab:latency}
\begin{tabular}{lr}
\toprule
\textbf{Metric} & \textbf{Value} \\
\midrule
Mean site build time            & \compSiteMeanBuildMs\,ms \\
Max site build time             & \compSiteMaxBuildMs\,ms \\
Encode-for-solver (mean)        & \compProfEncodeForSolverMeanMs\,ms \\
Formal core site (mean)         & \compProfFormalCoreSiteMeanMs\,ms \\
Orchestrate verification (mean) & \compProfOrchestrateVerificationMeanMs\,ms \\
Specification satisfaction (mean) & \compProfSpecificationSatisfactionMeanMs\,ms \\
Mean end-to-end time            & \compTimeMean \\
\bottomrule
\end{tabular}
\end{table}

The solver encoding dominates per-site overhead at
\compProfEncodeForSolverMeanMs\,ms, while all other federation stages
contribute sub-millisecond overhead.  The end-to-end pipeline
averages \compTimeMean{} per program, with site construction
accounting for a negligible fraction.

\subsection{Conflict statistics}

Of the \numcases{} cases, 23 generated at least one conflict (7.7\%).
Table~\ref{tab:conflicts} breaks these down by kind.

\begin{table}[h]
\centering
\caption{Conflict resolution: qualitative distribution by kind.}
\label{tab:conflicts}
\begin{tabular}{lll}
\toprule
\textbf{ConflictKind} & \textbf{Resolution policy} & \textbf{Frequency} \\
\midrule
TrustConflict     & Automatic (trust lattice join) & Most common \\
AuthorityConflict & Automatic (authority hierarchy) & Common \\
SpecMismatch      & Manual review required & Rare \\
CircularDep       & Manual review required & Rare \\
\bottomrule
\end{tabular}
\end{table}

Trust and authority conflicts are resolved automatically via the
algebraic operations of \S\ref{sec:algebra}.  Specification mismatches
and circular dependencies require manual intervention but are the
least frequent in practice.  Across the benchmark,
\compTrustSolverdischarged{} trust obligations were discharged
(\compTrustSolverdischargedPct{}) with \compTrustUnverified{}
remaining unverified.

% ══════════════════════════════════════════════════════════════════════
\section{Conclusion}
\label{sec:conclusion}
% ══════════════════════════════════════════════════════════════════════

We have presented Pack Federation, a framework for distributing
\JuGeo{} verification across organisational boundaries via typed
bridges between pack specifications.  The framework's ten components
(discovery, registry, composition, verification, evidence combining,
authority management, maintenance, serialisation, and statistics) are
implemented in \texttt{src/jugeo/packs/} and evaluated on a 12-pack
corpus.

The Federation Consistency Theorem establishes that local verification
combined with structural bridge validity is sufficient for global
consistency.  The proof is direct and machine-checked; its simplicity
is intentional---it reflects a deliberate design choice to make the
correctness argument obvious rather than clever.

Future work will address (i) dynamic pack admission (adding a new pack
to a live federation without re-verifying existing bridges), (ii)
cryptographic authority signatures for cross-organisation trust, and
(iii) extension of bridge composition to handle obligation transfer via
\TRANSPORT{} and \ESCALATE{} moves.

\begin{thebibliography}{99}

\bibitem{jugeo01}
H.~Young.
\newblock Judgment Geometry~I: Semantic Sites for Program Verification.
\newblock \textit{JuGeo Series}, Paper~1, 2024.

\bibitem{jugeo02}
H.~Young.
\newblock Judgment Geometry~II: Grothendieck Topologies on Code Coordinates.
\newblock \textit{JuGeo Series}, Paper~2, 2024.

\bibitem{jugeo10}
H.~Young.
\newblock Trust Algebras for Verification Evidence.
\newblock \textit{JuGeo Series}, Paper~10, 2024.

\bibitem{jugeo27}
H.~Young.
\newblock Certificate Chains and Provenance Tracking.
\newblock \textit{JuGeo Series}, Paper~27, 2024.

\bibitem{jugeo36}
H.~Young.
\newblock Ablation Methodology for Judgment Geometry.
\newblock \textit{JuGeo Series}, Paper~36, 2025.

\bibitem{leroy}
X.~Leroy.
\newblock Formal verification of a realistic compiler.
\newblock \textit{Communications of the ACM}, 52(7):107--115, 2009.

\bibitem{z3}
L.~de~Moura and N.~Bj{\o}rner.
\newblock {Z3}: An efficient {SMT} solver.
\newblock In \textit{TACAS}, LNCS 4963, pp.~337--340, 2008.

\bibitem{lean4}
L.~de~Moura and S.~Ullrich.
\newblock The {Lean 4} theorem prover and programming language.
\newblock In \textit{CADE-28}, LNCS 12699, pp.~625--635, 2021.

\bibitem{moduverif}
N.~Beckman, A.~Nori, S.~Rajamani, and R.~Simmons.
\newblock Proofs from tests.
\newblock In \textit{ISSTA}, pp.~3--14, 2008.

\bibitem{appel-pccc}
A.~W.~Appel and A.~Felty.
\newblock A semantic model of types and machine instructions for proof-carrying
  code.
\newblock In \textit{POPL}, pp.~243--253, 2000.

\bibitem{stp}
V.~Ganesh and D.~L.~Dill.
\newblock A decision procedure for bit-vectors and arrays.
\newblock In \textit{CAV}, LNCS 4590, pp.~519--531, 2007.

\bibitem{grothendieck-sga4}
M.~Artin, A.~Grothendieck, and J.-L.~Verdier.
\newblock \textit{Th\'{e}orie des Topos et Cohomologie \'Etale des Sch\'emas
  (SGA~4)}.
\newblock Lecture Notes in Mathematics 269--270--305.
  Springer, 1972--1973.

\bibitem{maclane-moerdijk}
S.~Mac~Lane and I.~Moerdijk.
\newblock \textit{Sheaves in Geometry and Logic}.
\newblock Springer, 1994.

\bibitem{nipkow-isabelle}
T.~Nipkow, L.~C.~Paulson, and M.~Wenzel.
\newblock \textit{Isabelle/HOL: A Proof Assistant for Higher-Order Logic}.
\newblock LNCS 2283. Springer, 2002.

\bibitem{pierce-tpl}
B.~C.~Pierce.
\newblock \textit{Types and Programming Languages}.
\newblock MIT Press, 2002.

\end{thebibliography}

\end{document}
